\section{یادگیری مشارکتی با توزیع طبقه‌بندی‌شدۀ داده‌ها}
\label{sec:federated}
در معماری مشارکتی، گزارش تعداد کلاینت‌ها به‌تنهایی برای توصیف آزمایش کافی نیست. تعداد نمونه‌های هر کلاینت، نسبت کلاس‌ها، استقلال داده‌های محلی، محل یادگیری بازنمایی متن و روش تجمیع پیش‌بینی‌ها باید مشخص باشند. در این پژوهش، این موارد بر اساس پیاده‌سازی و خروجی‌های ثبت‌شده در نوت‌بوک \en{federated\_xgboost\_10\_20\_50\_100.ipynb} تعریف می‌شوند \cite{notebook-federated}. تمام اعداد و نمودارهای این بخش از همان اجرای ذخیره‌شده استخراج شده‌اند؛ اعداد جدول‌ها تا چهار رقم اعشار گرد شده‌اند و نباید به‌عنوان دقت بیشتر از خروجی اصلی تفسیر شوند.

\subsection{مجموعه‌داده و جداسازی ارزیابی}
دادۀ ورودی، مجموعۀ \en{Suicide Watch} با ۲۳۲٬۰۷۴ ردیف و دو ستون اصلی متن و برچسب است. پس از حذف ۹۳ متن خام تکراری، ۲۳۱٬۹۸۱ ردیف باقی می‌ماند. ابتدا ۲۰ درصد داده به آزمون اختصاص می‌یابد؛ سپس ۲۰ درصد بخش توسعه برای اعتبارسنجی کنار گذاشته می‌شود. بنابراین سهم‌های اسمی آموزش، اعتبارسنجی و آزمون به‌ترتیب ۶۴، ۱۶ و ۲۰ درصد است. بذر تصادفی تقسیم برابر ۴۲ است و تقسیم‌ها با حفظ نسبت کلاس انجام می‌شوند.

پس از پاکسازی، متن‌های خالی، متن‌های تکراریِ درون هر بخش، متن‌های دارای برچسب متعارض و هم‌پوشانی با بخش‌های قبلی حذف می‌شوند. در نتیجه، تعدادهای نهایی الزاماً برابر حاصل‌ضرب سادۀ نسبت‌های اسمی در تعداد اولیه نیستند. جدول~\ref{tab:fl-split} تعدادهای واقعی پس از این کنترل‌ها را نشان می‌دهد. در این آزمایش، داده‌ها انگلیسی هستند؛ ارزیابی چندزبانه یا تعمیم بالینی مستقل انجام نشده است.

\begin{table}[htbp]
\centering
\caption{توزیع واقعی داده پس از حذف تکرار و پاکسازی در آزمایش مشارکتی}
\label{tab:fl-split}
\begin{tabular}{lrrrr}
\toprule
بخش & غیرخودکشی & خودکشی & مجموع & سهم کلاس مثبت \\
\midrule
آموزش & 74133 & 74247 & 148380 & 0.500384 \\
اعتبارسنجی & 18525 & 18561 & 37086 & 0.500485 \\
آزمون & 23145 & 23197 & 46342 & 0.500561 \\
\bottomrule
\end{tabular}
\end{table}

مجموعه‌های اعتبارسنجی و آزمون در تمام سناریوها مشترک و از دادۀ آموزش کلاینت‌ها جدا هستند. کنترل شناسۀ ردیف و کلید متن خام و پاکسازی‌شده مانع حضور نمونۀ مشترک در آموزش و ارزیابی می‌شود. برچسب‌های اعتبارسنجی و آزمون در اختیار سرور ارزیابی باقی می‌مانند و وارد تابع آموزش کلاینت نمی‌شوند. با این حال، محتوای پرسش ارزیابی برای محاسبۀ پیش‌بینی به مدل‌های محلی داده می‌شود؛ این ویژگی با محرمانه ماندن برچسب‌ها و دادۀ آموزش تفاوت دارد.

\begin{figure}[htbp]
\centering
\includegraphics[width=0.95\textwidth]{federated/cell-09-figure-0.png}
\caption{توزیع کلاس‌ها در بخش‌های مشترک و طول متن‌های آموزشی؛ خروجی نوت‌بوک مشارکتی}
\label{fig:fl-split}
\end{figure}

\subsection{تعریف دقیق تخصیص داده به کلاینت‌ها}
اگر $D_{\mathrm{tr}}$ مجموعۀ آموزش با اندازۀ $N=148380$ باشد، برای هر سناریوی $K\in\{10,20,50,100\}$، مجموعه‌های محلی $D_1,\ldots,D_K$ به‌گونه‌ای ساخته می‌شوند که
\begin{equation}
\bigcup_{i=1}^{K}D_i=D_{\mathrm{tr}},\qquad
D_i\cap D_j=\varnothing\quad(i\ne j),\qquad
\sum_{i=1}^{K}n_i=N,
\label{eq:fl-partition}
\end{equation}
که $n_i=|D_i|$ است. در تابع \en{stratified\_client\_indices}، اندیس‌های هر کلاس به‌طور جداگانه با مولد تصادفیِ بذر ۴۲ درهم‌ریخته می‌شوند و سپس با \en{numpy.array\_split} به $K$ بخش تقسیم می‌شوند. هر کلاینت یک بخش از هر کلاس را دریافت می‌کند و در پایان ترتیب ردیف‌های محلی نیز درهم‌ریخته می‌شود.

برای کلاس $c$ با $N_c$ نمونه، تعداد محلی از رابطۀ زیر پیروی می‌کند:
\begin{equation}
n_{i,c}\in\left\{\left\lfloor\frac{N_c}{K}\right\rfloor,
\left\lceil\frac{N_c}{K}\right\rceil\right\},\qquad
\pi_i=\frac{n_{i,1}}{n_i}\simeq\pi=\frac{74247}{148380}.
\label{eq:fl-stratified}
\end{equation}
از آنجا که گرد کردن برای هر کلاس جداگانه انجام می‌شود، اختلاف تعداد کل نمونه‌ها بین دو کلاینت می‌تواند دو نمونه باشد. این روش، یک تقسیم طبقه‌بندی‌شده و تقریباً هم‌توزیع است؛ یک سناریوی \en{non-IID}، تقسیم بر اساس کاربر، تقسیم زمانی یا شبیه‌سازی تفاوت واقعی بیمارستان‌ها محسوب نمی‌شود. شباهت نسبت کلاس‌ها نیز به معنای یکسان بودن واژگان و موضوع متن‌های تمام کلاینت‌ها نیست.

\begin{table}[htbp]
\centering
\caption{اندازه و نسبت کلاس مثبت در دادۀ محلی؛ تمام ۱۴۸٬۳۸۰ نمونه در هر سناریو دقیقاً یک‌بار استفاده می‌شوند}
\label{tab:fl-partition}
\small
\begin{tabular}{rrrrrr}
\toprule
کلاینت & کمینۀ ردیف & میانگین ردیف & بیشینۀ ردیف & کمینۀ $\pi_i$ & بیشینۀ $\pi_i$ \\
\midrule
10 & 14837 & 14838.0 & 14839 & 0.500371 & 0.500404 \\
20 & 7418 & 7419.0 & 7420 & 0.500337 & 0.500404 \\
50 & 2966 & 2967.6 & 2968 & 0.500337 & 0.500506 \\
100 & 1483 & 1483.8 & 1485 & 0.500337 & 0.500674 \\
\bottomrule
\end{tabular}
\end{table}

\subsection{آموزش مستقل بازنمایی و مدل محلی}
هر کلاینت مدل \en{Word2Vec} خود را فقط بر متن‌های آموزشی محلی یاد می‌گیرد. بازنمایی، از نوع \en{Skip-Gram} با ابعاد ۱۰۰، پنجرۀ ۵، حداقل فراوانی ۲، پنج دورۀ آموزش و یک کارگر است. بردار سند از میانگین بردار واژه‌های شناخته‌شده به دست می‌آید. واژه‌ای که فقط در اعتبارسنجی یا آزمون وجود داشته باشد وارد واژگان یادگرفته‌شده نمی‌شود. اگر تمام واژه‌های سند ناشناخته باشند، نمایش آن بردار صفر خواهد بود.

هر کلاینت یک \en{XGBoost} مستقل با ۱۰۰ درخت، عمق بیشینۀ ۵، نرخ یادگیری ۰٫۱ و روش ساخت درخت \en{hist} آموزش می‌دهد. تابع هدف \en{binary:logistic} و معیار داخلی \en{logloss} است. نسبت نمونه‌گیری سطری و ستونی هر دو یک، مقدار \en{min\_child\_weight} برابر یک و جریمۀ $L_2$ برابر یک است. هر مدل یک رشتۀ پردازشی دارد؛ اجرای ذخیره‌شده حداکثر چهار کلاینت را هم‌زمان شبیه‌سازی کرده است. بذر مدل کلاینت $i$، با شماره‌گذاری از یک، برابر $42+1009i$ است.

این پیاده‌سازی وزن درخت‌ها یا مختصات بردارهای \en{Word2Vec} را میانگین نمی‌گیرد. از آنجا که فضای بازنمایی هر کلاینت مستقل آموخته می‌شود، مختصات آن فضاها لزوماً هم‌تراز نیستند. اشتراک پیش‌بینی احتمالی، امکان ترکیب خروجی‌ها را بدون فرض هم‌ترازی این فضاها فراهم می‌کند. بنابراین روش حاضر، در سطح پیاده‌سازی، یک مجموعه‌مدل توزیع‌شده با تجمیع احتمال است و نباید با الگوریتم آموزش گرادیانی چنددوره‌ای \en{FedAvg} یکسان معرفی شود.

\subsection{تجمیع خروجی و انتخاب آستانه}
تمام کلاینت‌ها باید به پرسش‌های مشترک و با ترتیب شناسۀ یکسان پاسخ دهند. برای متن $x$، کلاینت $i$ احتمال کلاس مثبت $p_i(x)$ را می‌فرستد و سرور با وزن تعداد نمونۀ آموزشی آن کلاینت محاسبه می‌کند:
\begin{equation}
\alpha_i=\frac{n_i}{\sum_{j=1}^{K}n_j},\qquad
\bar p(x)=\sum_{i=1}^{K}\alpha_i p_i(x),\qquad
\widehat y(x)=\mathbf{1}\{\bar p(x)\ge\tau\}.
\label{eq:fl-aggregation}
\end{equation}
وزن‌ها از اندازۀ دادۀ محلی تعیین می‌شوند؛ از دقت آزمون یا برچسب آزمون برای تنظیم وزن استفاده نمی‌شود. نزدیک بودن اندازۀ کلاینت‌ها باعث می‌شود وزن‌ها تقریباً برابر باشند، اما فرمول واقعی همچنان از $n_i$ استفاده می‌کند.

آستانه از مجموعۀ $\{0,0.01,\ldots,1\}$ و تنها بر اساس بیشترین $F_1$ کلاس مثبت در اعتبارسنجی انتخاب می‌شود. در تساوی، دقت مثبت بیشتر، نزدیکی بیشتر به ۰٫۵ و سپس آستانۀ کوچک‌تر ترجیح دارند. پس از تثبیت آستانه، معیارهای آزمون محاسبه می‌شوند. برای تفکیک اثر تغییر تعداد کلاینت از اثر تنظیم آستانه، نتایج با آستانۀ مشترک ۰٫۵ نیز گزارش می‌شوند. سناریوی پیشنهادی نیز با $F_1$ اعتبارسنجی انتخاب می‌شود؛ بیشترین دقت آزمون مبنای انتخاب نیست.

\begin{table}[htbp]
\centering
\caption{نتایج آزمون با آستانۀ منتخب اعتبارسنجی؛ معیارها در مقیاس صفر تا یک}
\label{tab:fl-selected}
\small
\begin{tabular}{rrrrrrr}
\toprule
کلاینت & آستانه & دقت کلی & دقت مثبت & بازیابی مثبت & $F_1$ مثبت & $F_1$ کلان \\
\midrule
10 & 0.49 & 0.9158 & 0.9107 & 0.9222 & 0.9164 & 0.9158 \\
20 & 0.52 & 0.9093 & 0.9091 & 0.9096 & 0.9094 & 0.9093 \\
50 & 0.54 & 0.8930 & 0.8963 & 0.8891 & 0.8927 & 0.8930 \\
100 & 0.50 & 0.8700 & 0.8578 & 0.8874 & 0.8724 & 0.8700 \\
\bottomrule
\end{tabular}
\end{table}

\begin{table}[htbp]
\centering
\caption{معیارهای مستقل از آستانه، تعداد خطاها و زمان اجرای سناریوهای مشارکتی}
\label{tab:fl-ranking}
\small
\begin{tabular}{rrrrrr}
\toprule
کلاینت & \en{ROC AUC} & \en{AP} & منفی کاذب & مثبت کاذب & زمان دیواری (ثانیه) \\
\midrule
10 & 0.9727 & 0.9722 & 1805 & 2097 & 107.6356 \\
20 & 0.9686 & 0.9680 & 2096 & 2109 & 91.1462 \\
50 & 0.9581 & 0.9582 & 2572 & 2386 & 98.4640 \\
100 & 0.9435 & 0.9450 & 2612 & 3412 & 101.8770 \\
\bottomrule
\end{tabular}
\end{table}

\begin{table}[htbp]
\centering
\caption{نتایج آزمون با آستانۀ ثابت ۰٫۵؛ بدون تنظیم جداگانۀ آستانه برای هر سناریو}
\label{tab:fl-fixed}
\begin{tabular}{rrrrrr}
\toprule
کلاینت & دقت کلی & دقت مثبت & بازیابی مثبت & $F_1$ مثبت & $F_1$ کلان \\
\midrule
10 & 0.9163 & 0.9135 & 0.9199 & 0.9167 & 0.9163 \\
20 & 0.9092 & 0.9050 & 0.9146 & 0.9098 & 0.9092 \\
50 & 0.8923 & 0.8857 & 0.9010 & 0.8933 & 0.8922 \\
100 & 0.8700 & 0.8578 & 0.8874 & 0.8724 & 0.8700 \\
\bottomrule
\end{tabular}
\end{table}

\subsection{گزارش مستقل چهار سناریو}
\subsubsection{سناریوی ۱۰ کلاینت}
هر کلاینت به‌طور میانگین ۱۴٬۸۳۸ نمونه دارد. $F_1$ مثبت اعتبارسنجی در آستانۀ منتخب ۰٫۴۹ برابر ۰٫۹۱۹۸ است و این مقدار، سناریوی ۱۰ کلاینت را بر اساس قاعدۀ اعتبارسنجی به گزینۀ پیشنهادی تبدیل می‌کند. در آزمون، دقت کلی ۰٫۹۱۵۸ و بازیابی مثبت ۰٫۹۲۲۲ به دست آمده است. ماتریس سردرگمی با ترتیب سطر و ستون «غیرخودکشی، خودکشی» برابر است با
\begin{equation}
\begin{pmatrix}21048&2097\\1805&21392\end{pmatrix}.
\end{equation}
افزایش بازیابی نسبت به آستانۀ ۰٫۵ با افزایش مثبت کاذب همراه است. این مبادله نباید به‌عنوان بهبود هم‌زمان تمام معیارها بیان شود. میانگین اندازۀ واژگان محلی ۱۷٬۶۱۸٫۲ واژه و میانگین پوشش توکن‌های آزمون حدود ۰٫۹۸۶۰ است.
\begin{figure}[p]
\centering\includegraphics[width=\textwidth]{federated/cell-18-figure-0.png}
\caption{گزارش مستقل ۱۰ کلاینت: ماتریس سردرگمی و منحنی‌های ارزیابی و آستانه}
\label{fig:fl10}
\end{figure}

\subsubsection{سناریوی ۲۰ کلاینت}
میانگین دادۀ محلی به ۷٬۴۱۹ ردیف کاهش می‌یابد. آستانۀ منتخب ۰٫۵۲ و $F_1$ مثبت اعتبارسنجی ۰٫۹۱۳۴ است. دقت آزمون ۰٫۹۰۹۳ و $F_1$ مثبت آزمون ۰٫۹۰۹۴ است. در این سناریو، ۲٬۰۹۶ منفی کاذب و ۲٬۱۰۹ مثبت کاذب ثبت شده‌اند. نسبت به سناریوی ۱۰ کلاینت، کاهش حجم دادۀ محلی با کاهش اندازۀ واژگان محلی به میانگین ۱۲٬۴۳۰٫۲۵ همراه است. زمان دیواری ۹۱٫۱۴۶۲ ثانیه بوده، اما زمان یک اجرای منفرد برای استنتاج قانون قطعی سرعت کافی نیست.
\begin{figure}[p]
\centering\includegraphics[width=\textwidth]{federated/cell-20-figure-0.png}
\caption{گزارش مستقل ۲۰ کلاینت با آستانۀ منتخب ۰٫۵۲}
\label{fig:fl20}
\end{figure}

\subsubsection{سناریوی ۵۰ کلاینت}
میانگین دادۀ محلی ۲٬۹۶۷٫۶ ردیف و آستانۀ منتخب ۰٫۵۴ است. $F_1$ مثبت اعتبارسنجی ۰٫۸۹۴۵ و دقت آزمون ۰٫۸۹۳۰ به دست آمده‌اند. در آزمون، بازیابی مثبت ۰٫۸۸۹۱ و $F_1$ مثبت ۰٫۸۹۲۷ است. تعداد منفی‌های کاذب به ۲٬۵۷۲ می‌رسد. میانگین اندازۀ واژگان محلی ۷٬۶۵۴٫۵۶ است. افزایش آستانه نسبت به ۰٫۵ دقت مثبت را افزایش داده، اما بازیابی مثبت و $F_1$ مثبت آزمون را اندکی کاهش داده است؛ انتخاب آستانه همچنان پیش از دیدن نتایج آزمون انجام شده است.
\begin{figure}[p]
\centering\includegraphics[width=\textwidth]{federated/cell-22-figure-0.png}
\caption{گزارش مستقل ۵۰ کلاینت با آستانۀ منتخب ۰٫۵۴}
\label{fig:fl50}
\end{figure}

\subsubsection{سناریوی ۱۰۰ کلاینت}
هر کلاینت به‌طور میانگین ۱٬۴۸۳٫۸ نمونه دارد. آستانۀ اعتبارسنجی همان ۰٫۵ است و بنابراین دو تنظیم آستانه نتیجۀ یکسانی دارند. $F_1$ مثبت اعتبارسنجی ۰٫۸۷۳۷ و دقت آزمون ۰٫۸۷۰۰ است. بازیابی مثبت ۰٫۸۸۷۴ و دقت مثبت ۰٫۸۵۷۸ است. تعداد مثبت کاذب ۳٬۴۱۲ و تعداد منفی کاذب ۲٬۶۱۲ است. میانگین اندازۀ واژگان محلی به ۵٬۱۹۲٫۱۳ کاهش یافته است. در نتیجه، حفظ نسبت تقریباً یکسان کلاس‌ها به‌تنهایی برای حفظ کیفیت مدل‌های محلی در اندازه‌های کوچک کافی نبوده است.
\begin{figure}[p]
\centering\includegraphics[width=\textwidth]{federated/cell-24-figure-0.png}
\caption{گزارش مستقل ۱۰۰ کلاینت با آستانۀ ۰٫۵}
\label{fig:fl100}
\end{figure}

\subsection{تحلیل مقایسه‌ای و حدود اعتبار نتایج}
با ثابت ماندن کل دادۀ آموزش، افزایش تعداد کلاینت‌ها، حجم دادۀ محلی هر مدل و اندازۀ واژگان محلی آن را کاهش می‌دهد. در اجرای ثبت‌شده، دقت آزمون با آستانۀ منتخب از ۰٫۹۱۵۸ در ۱۰ کلاینت به ۰٫۸۷۰۰ در ۱۰۰ کلاینت می‌رسد؛ افتی برابر ۴٫۵۸ واحد درصد. $F_1$ مثبت نیز از ۰٫۹۱۶۴ به ۰٫۸۷۲۴ کاهش می‌یابد. افت \en{ROC AUC} از ۰٫۹۷۲۷ به ۰٫۹۴۳۵ نشان می‌دهد که تفاوت صرفاً ناشی از آستانه نیست. با این حال، این مشاهده به همین طراحی، مجموعه‌داده، پارامترها و بذر محدود است و یک حکم کلی دربارۀ تمام روش‌های یادگیری مشارکتی محسوب نمی‌شود.

\begin{figure}[htbp]
\centering\includegraphics[width=\textwidth]{federated/cell-28-figure-0.png}
\caption{مقایسۀ معیارهای آزمون سناریوهای مشارکتی بر اساس خروجی ثبت‌شده}
\label{fig:fl-comparison}
\end{figure}
\begin{figure}[htbp]
\centering\includegraphics[width=0.95\textwidth]{federated/cell-28-figure-1.png}
\caption{مقایسۀ منحنی‌های \en{ROC} و دقت مثبت\lr{--}بازیابی آزمون مشترک در چهار سناریو}
\label{fig:fl-coverage}
\end{figure}
\begin{figure}[htbp]
\centering\includegraphics[width=0.95\textwidth]{federated/cell-28-figure-2.png}
\caption{مقایسۀ ماتریس‌های سردرگمی چهار سناریو با آستانه‌های منتخب اعتبارسنجی}
\label{fig:fl-diagnostics}
\end{figure}
\FloatBarrier

برای هر سناریو فقط یک بذر، یعنی ۴۲، اجرا شده است. بنابراین انحراف معیار بین اجراها، فاصلۀ اطمینان تجربی و آزمون معنی‌داری بین سناریوها در دسترس نیست. گزارش عددی دقت نباید با قطعیت آماری یا ارزیابی بالینی اشتباه شود. این داده‌ها از متن‌های شبکۀ اجتماعی به دست آمده‌اند و برچسب متنی لزوماً معادل تشخیص بالینی خطر فرد نیست.

مقایسه با مدل متمرکز نیز باید با توجه به تنظیم‌ها انجام شود. مدل متمرکزِ منتخب، ۱۵۰ درخت، عمق ۱۰ و نرخ یادگیری ۰٫۲ دارد؛ مدل‌های محلی این بخش ۱۰۰ درخت، عمق ۵ و نرخ یادگیری ۰٫۱ دارند. بازنمایی متمرکز بر کل آموزش و بازنمایی‌های محلی بر بخش‌های جدا یاد گرفته می‌شوند. بنابراین اختلاف میان دقت \BestTestAccuracyPercent\ درصدِ مدل متمرکز و ۹۱٫۵۸ درصدِ سناریوی ۱۰ کلاینت، اثر خالصِ توزیع داده یا رمزگذاری را اندازه نمی‌گیرد. برای تفکیک علت‌ها به آزمایش کنترل‌شده با تنظیم‌های همسان نیاز است.

\subsection{ارتباط نتیجۀ مشارکتی با لایۀ رمزگذاری همریخت}
تجمیع احتمال در رابطۀ~\eqref{eq:fl-aggregation} در نوت‌بوک با اعداد اعشاری و به‌صورت آشکار انجام شده است. فصل‌های بعد، سازوکار رمزگذاری و احراز اصالت همین نوع تجمیع عددی را تشریح می‌کنند؛ خروجی‌های این نوت‌بوک به‌تنهایی آزمایش یک سامانۀ سرتاسری رمزگذاری‌شده نیستند. زمان‌های جدول~\ref{tab:fl-ranking} نیز زمان آموزش و پیش‌بینی مشارکتی‌اند و نباید با زمان رمزنگاری جمع شوند یا به‌عنوان زمان امنِ سرتاسری گزارش شوند.

در صورت استفاده از وزن صحیح $w_i=n_i$ و مقیاس $S=10^4$، برای جلوگیری از بازگشت پیمانه‌ای در \en{Benaloh} باید
\begin{equation}
r>S\sum_iw_i=10^4\times148380=1483800000
\label{eq:fl-benaloh-capacity}
\end{equation}
برقرار باشد. پارامتر $r=10000019$ در بنچمارک کوچک رمزنگاری برای این مجموع وزن کفایت نمی‌کند. اتصال عملی این دو آزمایش به انتخاب پارامتر پیام بزرگ‌تر، بررسی هزینۀ رمزگشایی و اندازه‌گیری اثر کوانتیزه‌سازی نیاز دارد. نرمال‌سازی وزن‌ها به اعداد صحیح کوچک‌تر نیز یک تقریب تازه ایجاد می‌کند و بدون ارزیابی نمی‌توان آن را معادل وزن‌دهی اصلی دانست. این تمایز، مرز میان نتایج یادگیری و نتایج لایۀ امنیت را روشن نگه می‌دارد.
\FloatBarrier
